\begin{center}
{
% POGING 1 - 1 rij met logo's en tekst%
%%%%%%%%%%%%%%%%%%%%%%%%%%%%%%%%%%%%%%%
\makeatletter
\vskip -140\p@
\begin{tabular}{p{0.2\textwidth}p{0.5\textwidth}p{0.2\textwidth}}
  \vspace{-2.5em} {\centering \includegraphics[width=0.2\textwidth]{GroupTEngineering_p2.eps} \par}
	& \vspace{-2.5em}
	{\small \centering 
	  Paper submitted within the scope of the Master's Thesis \\ Master of Industrial Sciences \\ GROUP T - Leuven Engineering College - 2012-2013 \par}
	& \vspace{-2.5em} {\centering \includegraphics[width=0.13\textwidth]{AssociatieLogo.eps} \par} %\\ \hline
\end{tabular} \par}
\vspace*{0.5cm} \centering{\emph{GROUP T does not guarantee error-free content of this paper.}}\\
{\vspace{5em} 
\makeatletter
\bf \huge{\@title} \par} \vskip 1.5em 
\makeatletter
{\large \@author}
\end{center}
\vspace*{5em}

\centerline{\bf ABSTRACT} \vspace*{0.2em}
\noindent The abstract should be viewed as a miniature version of the paper. Since potential readers should be able to make their decision on the personal relevance based
on the Abstract, the Abstract should clearly tell the reader what information he can expect
to find in the paper. The Abstract should (1) state the principal objectives and scope of the
investigation, (2) describe the methods employed, (3) summarize the results, and (4) state the
principal conclusions. Most of the Abstract will be written in the past tense, because it refers
to work done. The Abstract should never give any information or conclusion that is not stated
in the paper. The authors should always keep in mind, that the Abstract is the most frequently
read part of a paper. It should contain at most 200 words. Today, the Abstract is generally
printed in one paragraph.  Define all symbols used in the Abstract. Do not cite references in
the Abstract. Do not include or refer to tables and figures.
\\ \vspace*{2em}

\centerline{\bf Keywords} \vspace*{0.2em}
\noindent
About four key words or phrases in alphabetical order, separated by commas. Given the importance
of web-searching, it is essential that authors make every effort to ensure articles are found online,
quickly and accurately. The "key" to this is the appropriate use of keywords. Authors should know
the key phrases for their subject area. Reference to an established common indexing standard in a
particular discipline is a useful starting point (for instance ChemWeb).